\documentclass[aspectratio=169]{beamer}
\usepackage[utf8]{inputenc}
\usepackage[T1]{fontenc}
\usepackage[polish]{babel}
\usepackage{graphicx}
\usepackage{tikz}
\usetikzlibrary{shapes.geometric, arrows.meta, positioning}

\usetheme{Madrid}
\usecolortheme{default}

% Colors
\definecolor{plblue}{RGB}{0,83,159}
\setbeamercolor{structure}{fg=plblue}

% Remove navigation
\setbeamertemplate{navigation symbols}{}

% Title info
\title[RiskScoop AI]{RiskScoop AI}
\subtitle{System analizy ryzyka powodziowego\\z~wykorzystaniem agentów AI}
\author[Nikodem Wojtczak]{Nikodem Wojtczak\\Numer indeksu: 247823}
\institute{Promotor: dr~inż. Piotr Mielczarski}
\date[2026]{2026}

\begin{document}

% =============================================================================
% SLAJD 1: Strona tytułowa
% =============================================================================
\begin{frame}
    \titlepage
\end{frame}

% =============================================================================
% SLAJD 2: Problem i cel pracy
% =============================================================================
\begin{frame}{Problem i cel pracy}
\begin{block}{Problem badawczy}
Brak narzędzi umożliwiających analizę ryzyka powodziowego w~języku naturalnym dla niespecjalistów
\end{block}

\vspace{0.5cm}

\begin{block}{Cel pracy}
Zaprojektowanie i~implementacja systemu z~agentem AI do analizy ryzyka powodziowego na podstawie zapytań w~języku naturalnym
\end{block}

\vspace{0.5cm}

\textbf{Główne funkcje:}
\begin{itemize}
    \item Wyszukiwanie obiektów infrastruktury (szpitale, szkoły, budynki)
    \item Analiza zagrożenia powodziowego (Copernicus GloFAS)
    \item Wizualizacja wyników na interaktywnej mapie
    \item Klasyfikacja ryzyka (niskie/średnie/wysokie)
\end{itemize}
\end{frame}

% =============================================================================
% SLAJD 3: Architektura systemu
% =============================================================================
\begin{frame}{Architektura systemu}
\begin{center}
\begin{tikzpicture}[
    layer/.style={draw, rounded corners, minimum width=9cm, minimum height=1cm, align=center, font=\small},
    arrow/.style={-{Stealth[length=2mm]}, thick}
]

\node[font=\small] (user) at (0,0) {\textbf{Użytkownik}};

\node[layer, fill=blue!15, below=0.4cm of user] (frontend) {
    \textbf{Warstwa prezentacji}\\
    \footnotesize React 18 \quad Mapbox GL JS \quad ChatPanel \quad LayerPanel
};

\node[layer, fill=green!15, below=0.4cm of frontend] (agentos) {
    \textbf{Warstwa AgentOS}\\
    \footnotesize FastAPI \quad Gemini 2.0 Flash \quad 6 Tools \quad SQLite
};

\node[layer, fill=orange!15, below=0.4cm of agentos] (services) {
    \textbf{Warstwa serwisów}\\
    \footnotesize Nominatim \quad Overture \quad GloFAS \quad Layer
};

\node[layer, fill=red!15, below=0.4cm of services] (external) {
    \textbf{Dane zewnętrzne}\\
    \footnotesize OpenStreetMap \quad Overture Maps \quad Copernicus CDS \quad GeoJSON
};

\draw[arrow] (user) -- (frontend);
\draw[arrow] (frontend) -- (agentos);
\draw[arrow] (agentos) -- (services);
\draw[arrow] (services) -- (external);

\end{tikzpicture}
\end{center}
\end{frame}

% =============================================================================
% SLAJD 4: Narzędzia agenta
% =============================================================================
\begin{frame}{Narzędzia agenta}
\begin{center}
\textbf{6 narzędzi agenta:}
\vspace{0.5cm}

\begin{enumerate}
    \item \texttt{search\_division} -- wyszukiwanie lokalizacji
    \vspace{0.3cm}
    \item \texttt{get\_division} -- granice administracyjne
    \vspace{0.3cm}
    \item \texttt{get\_overture\_data} -- infrastruktura (MotherDuck)
    \vspace{0.3cm}
    \item \texttt{get\_flood\_forecast} -- dane powodziowe (GloFAS)
    \vspace{0.3cm}
    \item \texttt{intersect\_flood\_zones} -- analiza przecięć
    \vspace{0.3cm}
    \item \texttt{list\_layers} -- zarządzanie warstwami
\end{enumerate}
\end{center}
\end{frame}

% =============================================================================
% SLAJD 5: Framework i źródła danych
% =============================================================================
\begin{frame}{Framework i~źródła danych}
\begin{columns}[T]
\begin{column}{0.5\textwidth}
\textbf{Framework Agno AgentOS:}
\begin{itemize}
    \item Automatyczne generowanie FastAPI
    \item Przechowywanie sesji (SQLite)
    \item Zarządzanie historią konwersacji
\end{itemize}

\vspace{0.5cm}

\textbf{Model LLM:}
\begin{itemize}
    \item Google Gemini 2.0 Flash
    \item Rozumienie języka naturalnego
    \item Długi kontekst (1M tokenów)
\end{itemize}
\end{column}

\begin{column}{0.5\textwidth}
\textbf{Źródła danych:}
\begin{itemize}
    \item \textbf{OpenStreetMap}\\
    Nominatim + Overpass\\
    Granice administracyjne

    \vspace{0.3cm}

    \item \textbf{Overture Maps}\\
    MotherDuck API\\
    Infrastruktura

    \vspace{0.3cm}

    \item \textbf{Copernicus GloFAS}\\
    Climate Data Store\\
    Prognozy powodziowe
\end{itemize}
\end{column}
\end{columns}
\end{frame}

% =============================================================================
% SLAJD 6: Klasyfikacja ryzyka
% =============================================================================
\begin{frame}{Klasyfikacja ryzyka powodziowego}
\begin{center}
\begin{block}{Trójstopniowa skala oparta na danych GloFAS}
\vspace{0.3cm}
\begin{itemize}
    \item \textbf{Niskie} -- klasy głębokości 1--2 (0--1~m)
    \vspace{0.3cm}
    \item \textbf{Średnie} -- klasa głębokości 3 (1--2~m)
    \vspace{0.3cm}
    \item \textbf{Wysokie} -- klasy głębokości 4--5 (>2~m)
\end{itemize}
\vspace{0.3cm}
\end{block}
\end{center}
\end{frame}

% =============================================================================
% SLAJD 7: Wyniki testów
% =============================================================================
\begin{frame}{Wyniki testów}
\begin{columns}[T]
\begin{column}{0.5\textwidth}
\textbf{Testy funkcjonalne (20 scenariuszy):}
\begin{itemize}
    \item Wynik: \textbf{19/20 PASS (95\%)}
    \item Kategorie: healthcare, education, buildings, transportation
    \item Lokalizacje: 20 miast w~Europie Środkowej
    \item 1 błąd: Zurich (kanton zamiast miasta)
\end{itemize}

\vspace{0.3cm}

\textbf{Testy interpretacji NL (25 zapytań):}
\begin{itemize}
    \item Wynik: \textbf{21/25 (84\%)}
    \item Zapytania proste: 100\% (15/15)
    \item Zapytania złożone: 60\% (6/10)
    \item Benchmark: 76--87\% \textcolor{green}{$\checkmark$}
\end{itemize}
\end{column}

\begin{column}{0.5\textwidth}
\textbf{Testy wydajnościowe:}
\begin{itemize}
    \item Średni czas: \textbf{14.2~s}
    \item Min: 8.4~s, Max: 23.1~s
    \item 90\% $\leq$ 30~s \textcolor{green}{(RN-01 $\checkmark$)}
    \item Rozkład czasu:
    \begin{itemize}
        \item Overture: 60\%
        \item GloFAS: 20\%
        \item Analiza: 18\%
        \item OSM: 2\%
    \end{itemize}
\end{itemize}

\vspace{0.2cm}

\textbf{Przykłady (rozkład ryzyka):}
\begin{itemize}
    \item \textbf{Genewa:} 47 szpitali\\12 w~ryzyku: 7 N / 3 Ś / 2 W
    \item \textbf{Warszawa:} 892 szkoły\\67 w~ryzyku: 41 N / 18 Ś / 8 W
    \item \textbf{Kraków:} 15,847 budynków\\2,134 w~ryzyku: 1,456 N / 512 Ś / 166 W
\end{itemize}
\end{column}
\end{columns}
\end{frame}

% =============================================================================
% SLAJD 8: Wnioski i zakończenie
% =============================================================================
\begin{frame}{Wnioski}
\textbf{Osiągnięcia:}
\begin{itemize}
    \item Udana implementacja systemu analizy ryzyka powodziowego z~agentem AI
    \item Wysoka jakość interpretacji języka naturalnego (84\% vs 76--87\% benchmark)
    \item Pełna integracja z~3 źródłami danych (OSM, Overture Maps, GloFAS)
    \item 95\% testów funkcjonalnych zakończonych sukcesem
\end{itemize}

\vspace{1cm}

\begin{center}
\Large \textbf{Dziękuję za uwagę}
\end{center}
\end{frame}

\end{document}
